\section{Архитектура и исходные файлы}
\label{sec:architecture}

\subsection{Организация модулей}
Публичные сущности предметной области объявлены в \code{back/}. Их реализации
расположены в корне проекта. Координация эксперимента, учёт закупочных цен,
поставки и консольное представление вынесены в \code{internal/} и пространство
имён \code{model_detail}. Публичные функции из \code{back/simulation.h}
делегируют работу единственному \code{SimulationEngine}.

\begin{figure}[htbp]
\centering
\begin{tikzpicture}[
  node distance=7mm,
  box/.style={draw=accent,rounded corners=2pt,fill=accent!5,
              align=center,text width=43mm,minimum height=12mm,font=\small},
  arrow/.style={-{Latex[length=2mm]},draw=accent,thick}
]
\node[box] (entry) {\texttt{main.cpp}\\Командная строка и запуск};
\node[box,below=of entry] (public) {\texttt{back/simulation.h}\\Публичный контракт};
\node[box,below=of public] (engine) {\texttt{SimulationEngine}\\Фазы дня, очередь, результаты};
\node[box,below left=9mm and 5mm of engine] (domain) {\texttt{Warehouse}, \texttt{Order}\\Товар, выдача и цена};
\node[box,below right=9mm and 5mm of engine] (supply) {\texttt{SupplyManager}\\Заявки и поступления};
\node[box,below=28mm of engine] (context) {\texttt{ModelContext}\\Генератор, скидки, учёт партий};
\draw[arrow] (entry) -- (public);
\draw[arrow] (public) -- (engine);
\draw[arrow] (engine) -- (domain);
\draw[arrow] (engine) -- (supply);
\draw[arrow] (domain) -- (context);
\draw[arrow] (supply) -- (context);
\end{tikzpicture}
\caption{Основные связи модулей. Генерация данных и сериализация описаны в тексте.}
\end{figure}

\code{DataGenerator} использует случайный генератор из \code{ModelContext}
и текущий склад. \code{simulation_cli.cpp} читает параметры, представляет
каталог и результаты в консоли и записывает JSON. Его форматирование отделено
от правил работы склада и заказов.

\subsection{Карта файлов}
{\small
\begin{longtable}{@{}L{0.38\textwidth}L{0.55\textwidth}@{}}
\toprule
Файл или каталог & Назначение \\
\midrule
\endhead
\code{PharmDelivery.pro} & Список исходников, заголовков и параметры qmake. \\
\code{main.cpp} & Разбор аргументов, создание \code{QApplication}, дневной цикл,
итоговый вывод и обработка исключений. \\
\code{back/simulation.h} & Параметры эксперимента, структуры результатов,
публичные функции управления моделью. \\
\code{back/data.h} & Глобальные склад, каталог, покупатели и параметры $N,M,K$. \\
\code{back/medicine.h}, \code{back/customer.h} & Поля лекарства и покупателя. \\
\code{back/store.h}, \code{store.cpp} & Партия, склад, уценка, списание,
проверка необходимости пополнения. \\
\code{back/order.h}, \code{order.cpp} & Запрос покупателя, подбор партий,
фактическая выдача и расчёт стоимости. \\
\code{back/request.h} & Заявка поставщику. \\
\code{back/statistics.h}, \code{statistics.cpp} & Накопление выручки,
закупочных расходов и потерь. \\
\code{back/data_generator.h}, \code{data_generator.cpp} & Каталог,
покупатели, состав заказов и задержки поставок. \\
\code{internal/model_support.*} & Генератор, политика скидок, проверки
входных данных и дополнительные сведения о партиях. \\
\code{internal/supply_manager.*} & Очередь поставок и история поступлений. \\
\code{internal/simulation_engine.*} & Состояние эксперимента, порядок фаз,
очередь заказов, распределение доставок и баланс. \\
\code{internal/simulation_cli.*} & Аргументы, справка, консольные и JSON-отчёты. \\
\code{customer.cpp}, \code{medicine.cpp}, \code{request.cpp}, \code{data.cpp}
& В текущей версии пусты; соответствующие поля и inline-глобальные переменные
объявлены в заголовках. \\
\code{examples/initial-stock.csv} & Пример начальных партий для 15 лекарств. \\
\code{tests/} & Модульные, интеграционные и консольные проверки. \\
\code{docs/} & Исходники этой документации и инструкция её сборки. \\
\bottomrule
\end{longtable}
}

\subsection{Состояние и владение данными}
\code{back/data.h} содержит inline-переменные \code{warehouse},
\code{mas_med}, \code{type_med}, \code{group_med}, \code{customers},
\code{N}, \code{M}, \code{K} и \code{global_marcup}.
Написание \code{global_marcup} соответствует существующему контракту.

\code{SimulationEngine} владеет очередью \code{orders_}, списком подготовленных
покупок \code{tomorrow_}, статистикой и результатом. Внутренний \code{QueuedOrder}
хранит идентификатор, исходный день заказа и \code{Order}. Подготовленная покупка
содержит копию данных покупателя, запрошенные и реально выданные позиции,
цену и закупочную стоимость.

\code{ModelContext} хранит единый \code{std::mt19937}, текущий день,
задержку доставки, настройки скидок и дополнительный учёт складов по их адресам.
\code{WarehouseAccounting} содержит закупочные цены партий, признаки уценки,
стоимость и количество выделенных упаковок, а также потери от списания.
Это позволяет сохранить исходные поля \code{StoreBatch} без добавления
отдельного поля закупочной цены.

Пары сведений о партии сопоставляются по лекарству, сроку годности и продажной
базе. При интеграции следует использовать методы модели: произвольная замена
элементов \code{inventory} или изменение их цен может нарушить соответствие
между публичным складом и внутренним учётом.
